% Figure: ROC curve for a binary classifier, sweeping the decision
% threshold from strict to permissive. The curve rises from (0,0) to
% (1,1); the area under it (AUC) summarises ranking quality. The
% diagonal is the no-skill reference. Two operating points are marked:
% a high-specificity point and a high-sensitivity point, showing that
% the threshold, not the model, sets the trade between the two errors.
% Coordinates are a precomputed concave ROC for a moderately skilful
% classifier (AUC ~ 0.88).
\begin{figure}[htbp]
\centering
\begin{tikzpicture}
\begin{axis}[
    width=9cm, height=8.5cm,
    xmin=0, xmax=1,
    ymin=0, ymax=1,
    axis lines=left,
    xlabel={False-positive rate $1-\text{specificity}$},
    ylabel={True-positive rate (sensitivity)},
    every axis plot/.append style={very thick, mark=none},
    legend style={font=\small, draw=none, fill=none, at={(0.98,0.05)}, anchor=south east},
    xtick={0,0.2,0.4,0.6,0.8,1.0},
    ytick={0,0.2,0.4,0.6,0.8,1.0},
]
% No-skill diagonal.
\addplot[enggray!60, dashed] coordinates {(0,0) (1,1)};
\addlegendentry{no-skill reference}
% Concave ROC curve for a skilful classifier.
\addplot[engblue] coordinates {
 (0.00,0.00) (0.03,0.28) (0.07,0.48) (0.12,0.63) (0.18,0.74)
 (0.27,0.83) (0.38,0.90) (0.52,0.95) (0.70,0.98) (0.85,0.995)
 (1.00,1.00)};
\addlegendentry{fitted classifier (AUC $\approx 0.88$)}
% Operating points.
\addplot[only marks, engred, mark=*, mark size=2.5pt] coordinates {(0.07,0.48) (0.38,0.90)};
\node[anchor=west, font=\footnotesize, engred] at (axis cs:0.09,0.44) {strict threshold};
\node[anchor=west, font=\footnotesize, engred] at (axis cs:0.40,0.86) {permissive threshold};
\end{axis}
\end{tikzpicture}
\caption[Receiver operating characteristic of a classifier]{The
receiver operating characteristic traces the true-positive rate
against the false-positive rate as the decision threshold sweeps from
strict to permissive. Every point on the curve is one threshold on one
fitted model; moving along the curve trades sensitivity against
specificity without changing the model at all. The area under the
curve summarises the model's ability to rank a random positive above a
random negative and is threshold-free, which is why it compares models
where accuracy at a single cutoff would not. The two marked operating
points make the standing lesson concrete: the model sets the curve,
the engineer sets the point, and the choice of point follows the cost
of each error, not a default.}
\label{fig:vol02:ch14:roc-curve}
\end{figure}
